% ==============================================================================
% Plantilla Institucional de Trabajo Terminal — UPIIZ IPN
% Archivo: capitulos/10-trabajo-futuro.tex
% ==============================================================================
% NOTA: Este capítulo se incluye condicionalmente solo para TT II.
% En PROTOCOLO y TT I se omite automáticamente.
% ==============================================================================

\chapter{Trabajo Futuro}
\label{cap:trabajo-futuro}

En este capítulo se proponen las líneas de investigación, mejoras de diseño y extensiones tecnológicas recomendadas para dar continuidad al proyecto desarrollado:

\begin{enumerate}
    \item \textbf{Optimizaciones mecánicas y estructurales:} Rediseño de componentes para reducción de masa, selección de materiales avanzados o incorporación de mecanismos con mayor grado de libertad.
    \item \textbf{Mejoras en sensado y electrónica:} Incorporación de sensores de mayor resolución, implementación de etapas de potencia más eficientes o integración de buses de comunicación industrial.
    \item \textbf{Evolución de algoritmos y control:} Exploración de técnicas de control adaptativo, predictivo basado en modelos (MPC) o algoritmos de aprendizaje por refuerzo.
    \item \textbf{Escalabilidad y aplicaciones industriales:} Adaptación del prototipo para pruebas en entornos productivos reales o integración con plataformas de Internet de las Cosas (IoT) e Industria 4.0.
\end{enumerate}
